\documentclass[a4paper,12pt]{article}
\usepackage[T2A]{fontenc}
\usepackage[utf8]{inputenc}
\usepackage[russian]{babel}
\usepackage{amsmath,amssymb,mathtools}
\usepackage{helvet}
\renewcommand{\familydefault}{\sfdefault}
\usepackage{icomma}
\usepackage[margin=20mm,top=20mm,bottom=21mm]{geometry}
\usepackage{microtype}
\usepackage{tikz}
\usetikzlibrary{arrows.meta,positioning,shapes.geometric,calc}
\usepackage{booktabs,tabularx,longtable}
\usepackage{enumitem}
\usepackage{fancyhdr}
\usepackage{parskip}
\usepackage{xcolor}
\usepackage{hyperref}
\hypersetup{hidelinks}
\definecolor{accent}{HTML}{007E82}
\definecolor{darkaccent}{HTML}{00535A}
\definecolor{mainink}{HTML}{292F34}
\definecolor{grayink}{HTML}{5A6268}
\color{mainink}
\setlength{\parindent}{0pt}
\setlength{\parskip}{5.6pt}
\setlist[itemize]{leftmargin=18pt,topsep=3pt,itemsep=3pt}
\setlist[enumerate]{leftmargin=23pt,topsep=5pt,itemsep=8pt}
\fancyhf{}
\pagestyle{fancy}
\renewcommand{\headrulewidth}{0.25pt}
\fancyhead[L]{\footnotesize\color{grayink}Задачи на движение}
\fancyhead[R]{\footnotesize\color{grayink}Анна Баннова}
\fancyfoot[C]{\footnotesize\color{grayink}\thepage}
\setlength{\headheight}{15pt}
\newcommand{\key}[1]{\textcolor{darkaccent}{\textbf{#1}}}
\newcommand{\topic}[1]{\vspace{7pt}{\large\bfseries\color{darkaccent}#1}\par\vspace{2pt}}
\newcommand{\subtopic}[1]{\vspace{4pt}{\bfseries #1}\par}
\newcommand{\vst}{\vspace{4pt}}
\newcommand{\example}[1]{\vspace{3pt}{\bfseries\color{darkaccent}Пример #1.}\enspace}
\newcommand{\smallnote}[1]{{\small\color{grayink}#1}}
\begin{document}
\thispagestyle{empty}
\begin{center}
\vspace*{24mm}
{\Large\color{darkaccent}МАТЕМАТИКА}\par
\vspace{18mm}
{\Huge\color{mainink}Чек-лист}\par
\vspace{12mm}
{\LARGE\color{accent}Текстовые задачи на движение}\par
\vspace{3mm}
{\Large\color{accent}Схемы, таблицы и уравнения}\par
\vspace{16mm}
{\large Подготовка к ЕГЭ по профильной математике}\par
\vfill
{\large Анна Баннова}\par
\vspace{10mm}
{\large 6 ноября 2026 г.}\par
\vspace{17mm}
{\normalsize Лёвин Артём Александрович}\par
{\normalsize эксклюзивно для Анны Банновой}\par
\vspace{19mm}
\begin{tikzpicture}
\draw[accent, line width=0.85pt]
(0,0) .. controls (1.9,0.9) and (3.25,-0.85) .. (5.2,0)
.. controls (6.95,0.9) and (8.75,-0.85) .. (10.5,0);
\fill[darkaccent] (5.2,0) circle(1.65pt);
\end{tikzpicture}
\vspace*{7mm}
\end{center}
\clearpage

\topic{1. Три опоры при решении задачи}
Почти вся работа над текстовой задачей о движении укладывается в последовательность:
\begin{center}
{\color{darkaccent}\bfseries схема пути $\longrightarrow$ таблица $v$--$t$--$S$ $\longrightarrow$ уравнение.}
\end{center}
\begin{itemize}
\item \key{Схема} показывает, кто откуда вышел, в каком направлении двигался, где встретились участники и какие отрезки образуют путь.
\item \key{Таблица} переводит словесное описание в математические величины.
\item \key{Уравнение} выражает связь расстояний или времён, прямо следующую из условия и схемы.
\end{itemize}
\vst
\subtopic{Формулы движения с постоянной скоростью}
\[\boxed{S=vt,\qquad t=\frac Sv\ (v>0),\qquad v=\frac St\ (t>0).}\]
Обычно скорость измеряется в км/ч, время --- в часах, путь --- в километрах.
\begin{center}
\renewcommand{\arraystretch}{1.25}
\begin{tabularx}{\linewidth}{@{}lXXX@{}}
\toprule
Участник & Скорость, км/ч & Время, ч & Путь, км\\
\midrule
Первый & $v_1$ & $t_1$ & $v_1t_1$\\
Второй & $v_2$ & $t_2$ & $v_2t_2$\\
\bottomrule
\end{tabularx}
\end{center}
\key{Как выбирать неизвестную.} Если просят найти время, удобно положить $t=x$; если скорость автобуса --- $v_{\text{авт}}=x$; если расстояние до опушки --- $S=x$. При необходимости вводите дополнительную переменную $y$: она часто сокращается при составлении уравнения.

\topic{2. Единицы измерения}
\begin{center}
\renewcommand{\arraystretch}{1.25}
\begin{tabularx}{\linewidth}{@{}XX@{}}
\toprule
Из метров в километры & Из километров в метры\\
$S_{\text{км}}=S_{\text{м}}/1000$ & $S_{\text{м}}=1000S_{\text{км}}$\\
\midrule
Из минут в часы & Из часов в минуты\\
$t_{\text{ч}}=t_{\text{мин}}/60$ & $t_{\text{мин}}=60t_{\text{ч}}$\\
\bottomrule
\end{tabularx}
\end{center}
\example{1}Если расстояние между пешеходами составляет $200$ м, то
\[200\text{ м}=0{,}2\text{ км}.\]
Время $0{,}2$ ч составляет $0{,}2\cdot60=12$ мин. Разница в $8$ мин равна $8/60=2/15$ ч.

\smallnote{Перед составлением уравнения сверьте единицы: при скорости в км/ч все расстояния должны быть в километрах, а времена --- в часах.}
\clearpage

\topic{3. Догоняющее движение: разность пройденных путей}
\example{2}Грузовик выезжает из пункта $B$, легковой автомобиль --- из $A$ в том же направлении. Между пунктами $132$ км; автомобиль едет на $22$ км/ч быстрее. Через сколько часов он догонит грузовик?
\begin{center}
\begin{tikzpicture}[x=1cm,y=1cm]
\draw[grayink,line width=.8pt] (0,0)--(11,0);
\draw[accent,line width=1pt,-{Latex[length=2mm]}] (0,.35)--(10.8,.35);
\draw[grayink,line width=1pt,-{Latex[length=2mm]}] (3,.7)--(10.8,.7);
\foreach \x/\n in {0/A,3/B,10.8/C}{\fill[mainink] (\x,0) circle (1.5pt);\node[below] at (\x,0){$\n$};}
\draw[<->,darkaccent] (0,-.66)--node[fill=white,inner sep=2pt]{\small 132 км}(3,-.66);
\node[above,grayink] at (6.3,.75){\footnotesize путь грузовика};
\node[above,accent] at (5.4,.32){\footnotesize путь автомобиля};
\end{tikzpicture}
\end{center}
Введём $x$ --- время движения (ч), $y$ --- скорость грузовика (км/ч).
\begin{center}
\renewcommand{\arraystretch}{1.18}
\begin{tabular}{@{}lccc@{}}\toprule
& $v$ & $t$ & $S$\\ \midrule
Автомобиль & $y+22$ & $x$ & $x(y+22)$\\
Грузовик & $y$ & $x$ & $xy$\\ \bottomrule
\end{tabular}
\end{center}
На рисунке видно: путь автомобиля длиннее на $132$ км. Поэтому
\[x(y+22)=132+xy,\quad xy+22x=132+xy,\quad22x=132.\]
\[\boxed{x=6\text{ ч}.}\]
\key{Обобщение.} При одновременном движении в одном направлении, когда более быстрый догоняет медленного с начальным отставанием $d$,
\[\boxed{t=\frac{d}{v_{\text{быстр}}-v_{\text{медл}}}}\qquad(v_{\text{быстр}}>v_{\text{медл}}).\]
\vst
\topic{4. Два участника стартуют вместе}
\example{3}Скорость быстрого пешехода на $1$ км/ч больше. Через сколько минут расстояние между пешеходами станет $200$ м?
\[x(y+1)=xy+0{,}2\quad\Rightarrow\quad x=0{,}2\text{ ч}=12\text{ мин}.\]
\key{Ключевое различие:} сначала найдены часы, а ответ нужен в минутах.
\clearpage

\topic{5. Поворот и движение обратно}
\example{4}Два человека одновременно отправились от старта к опушке в $6$ км. Скорость быстрого $6{,}5$ км/ч, медленного $3{,}5$ км/ч. Быстрый дошёл до опушки и повернул назад. На каком расстоянии от опушки они встретились?
\begin{center}
\begin{tikzpicture}[x=1cm,y=1cm]
\draw[mainink,line width=1pt] (0,0)--(11,0);
\fill[mainink] (0,0) circle(1.6pt) node[below]{старт};
\fill[mainink] (11,0) circle(1.6pt) node[below]{опушка};
\fill[darkaccent] (7.5,0) circle(2pt) node[below=8pt]{встреча};
\draw[accent,line width=1.5pt,-{Latex[length=2mm]}] (0,.5)--(10.9,.5);
\draw[accent,line width=1.5pt,-{Latex[length=2mm]}] (11,.85)--(7.5,.85);
\draw[grayink,line width=1.5pt,-{Latex[length=2mm]}] (0,1.1)--(7.45,1.1);
\draw[<->,darkaccent] (0,-.82)--node[fill=white,inner sep=2pt]{\small $6$ км}(11,-.82);
\draw[<->,darkaccent] (7.5,-.36)--node[fill=white,inner sep=2pt]{\small $x$}(11,-.36);
\node[grayink,above] at (3.0,1.12){\footnotesize медленный};
\node[accent,above] at (9.1,.86){\footnotesize быстрый: туда и обратно};
\end{tikzpicture}
\end{center}
Пусть $x$ км --- расстояние \emph{от места встречи до опушки}. Из схемы:
\[S_{\text{быстр}}=6+x,\qquad S_{\text{медл}}=6-x,\qquad0\le x\le6.\]
\begin{center}
\renewcommand{\arraystretch}{1.2}
\begin{tabular}{@{}lccc@{}}\toprule
& $v$, км/ч & $S$, км & $t=S/v$, ч\\ \midrule
Быстрый & $6{,}5$ & $6+x$ & $\dfrac{6+x}{6{,}5}$\\[3pt]
Медленный & $3{,}5$ & $6-x$ & $\dfrac{6-x}{3{,}5}$\\[3pt]
\bottomrule\end{tabular}
\end{center}
Вышли одновременно и встретились в один момент, следовательно, времена равны:
\[\frac{6+x}{6{,}5}=\frac{6-x}{3{,}5}.\]
\[3{,}5(6+x)=6{,}5(6-x),\qquad 10x=18,\qquad x=1{,}8.\]
\[\boxed{x=1800\text{ м}.}\]
\key{Частая ошибка:} записать $6-x$ для быстрого. Он прошёл весь путь до опушки и \emph{ещё} $x$ км назад.
\clearpage

\topic{6. Неодновременный старт: разность времён}
\example{5}Грузовик проехал $20$ км от $A$ до $B$. Автобус выехал на $8$ мин позже, ехал на $5$ км/ч быстрее и прибыл вместе с грузовиком. Найти скорость автобуса.
\par
Пусть $x$ км/ч --- скорость автобуса; тогда скорость грузовика --- $x-5$ км/ч. \key{Ограничение: $x>5$}. Оба преодолели $20$ км.
\begin{center}
\renewcommand{\arraystretch}{1.2}
\begin{tabular}{@{}lccc@{}}\toprule
& Скорость, км/ч & Расстояние, км & Время, ч\\ \midrule
Грузовик & $x-5$ & $20$ & $\dfrac{20}{x-5}$\\[3pt]
Автобус & $x$ & $20$ & $\dfrac{20}{x}$\\[3pt]
\bottomrule\end{tabular}
\end{center}
Грузовик ехал дольше автобуса на $8$ мин, то есть на $2/15$ ч. Поэтому
\[\frac{20}{x-5}-\frac{20}{x}=\frac{2}{15}.\]
\subtopic{Шаг 1. Приведение к общему знаменателю}
\[\frac{20x-20(x-5)}{x(x-5)}=\frac{2}{15}\quad\Rightarrow\quad\frac{100}{x^2-5x}=\frac{2}{15}.\]
Общий знаменатель $x(x-5)$ допустим: $x>5$.
\subtopic{Шаг 2. Пропорция и квадратное уравнение}
\[2(x^2-5x)=1500\quad\Rightarrow\quad2x^2-10x-1500=0.\]
Все коэффициенты делятся на $2$, поэтому
\[x^2-5x-750=0.\]
\subtopic{Шаг 3. Дискриминант и разложение на множители}
\[D=(-5)^2-4\cdot1\cdot(-750)=3025=5^2\cdot11^2.\]
\[\sqrt{3025}=5\cdot11=55,\qquad x_{1,2}=\frac{5\pm55}{2}.\]
Получаем $x_1=30$ и $x_2=-25$. Второй корень не удовлетворяет условию $x>5$.
\[\boxed{v_{\text{автобуса}}=30\text{ км/ч}.}\]
\smallnote{Проверка: грузовик едет со скоростью 25 км/ч, его время $20/25=0{,}8$ ч; автобус едет $20/30=2/3$ ч. Разница $0{,}8-2/3=2/15$ ч = 8 мин.}
\clearpage

\topic{7. Приёмы вычислений, которые пригодились на занятии}
\subtopic{Общий знаменатель для двух дробей}
При $b\ne0$ и $d\ne0$:
\[\boxed{\frac ab-\frac cd=\frac{ad-bc}{bd}.}\]
Вычитаем \emph{числители с учётом скобок}; знаменатели перемножаем.
\[\frac{20}{x-5}-\frac{20}{x}=\frac{20x-20(x-5)}{x(x-5)}.\]
\subtopic{Перемножение крест-накрест}
Если $B\ne0$, $D\ne0$, то
\[\frac AB=\frac CD\quad\Longleftrightarrow\quad AD=BC.\]
Многочлен берите в скобки: $2(x^2-5x)=1500$. При раскрытии скобок двойка умножает \emph{оба} слагаемых.
\subtopic{Деление квадратного уравнения}
\[2x^2-10x-1500=0\quad\Longleftrightarrow\quad x^2-5x-750=0.\]
Делить на ненулевое число можно \emph{всё уравнение целиком}, то есть каждый его член. Если деление даёт дробные коэффициенты, оно также допустимо, но часто неудобно: делимость коэффициентов влияет на удобство, а не на право делить.
\subtopic{Извлечение корня из большого числа}
\[3025=5\cdot5\cdot11\cdot11=5^2\cdot11^2,\qquad\sqrt{3025}=55.\]
Для неотрицательных $a,b$ верно $\sqrt{ab}=\sqrt a\sqrt b$. Нельзя без проверки извлекать корень из суммы по слагаемым.
\topic{8. Самопроверка после решения}
\begin{itemize}
\item Действительно ли $x$ обозначает величину из вопроса?
\item Верно ли расположены участники и их направления на схеме?
\item Нет ли смешения метров с километрами, минут с часами?
\item На каком основании приравнены пути или времена?
\item Учтены ли ограничения: положительные скорости, ненулевые знаменатели, расстояние до границы маршрута?
\item Прошли ли найденные корни проверку по смыслу задачи?
\end{itemize}
\clearpage

\thispagestyle{fancy}
\begin{center}{\LARGE\color{darkaccent}\bfseries Алгоритм: от условия к ответу}\end{center}
\vspace{9mm}
\begin{center}
\begin{tikzpicture}[
 node distance=11mm,
 basic/.style={draw=darkaccent,rounded corners=2.5pt,align=center,text width=10.5cm,minimum height=11mm,inner sep=5pt,font=\small},
 terminal/.style={draw=darkaccent,ellipse,align=center,text width=6.2cm,minimum height=10mm,inner sep=4pt,font=\small},
 decision/.style={diamond,draw=darkaccent,aspect=2.7,align=center,text width=4.35cm,inner sep=2pt,font=\small},
 branch/.style={draw=grayink,rounded corners=2pt,align=center,text width=5.0cm,minimum height=16mm,inner sep=4pt,font=\small},
 arr/.style={-{Latex[length=2.2mm]},line width=.72pt,color=darkaccent}
]
\node[terminal] (start) {Прочитать условие и вопрос};
\node[basic,below=of start] (scheme) {Нарисовать схему: старт, направление, встреча, поворот, участки пути};
\node[basic,below=of scheme] (units) {Привести величины к согласованным единицам: км, ч, км/ч};
\node[basic,below=of units] (table) {Назначить $x$ искомой величине и заполнить таблицу $v$--$t$--$S$};
\node[decision,below=15mm of table] (choice) {Какая связь дана в условии?};
\node[branch,below left=13mm and -2mm of choice] (road) {\textbf{Связь путей}\\начальный отрыв, разница пройденного, возврат};
\node[branch,below right=13mm and -2mm of choice] (times) {\textbf{Связь времён}\\одновременный старт, встреча, задержка выезда};
\node[basic,below=29mm of choice] (equation) {Составить уравнение по выбранной связи};
\node[basic,below=of equation] (solve) {Решить уравнение: дроби, скобки, квадратное уравнение};
\node[terminal,below=of solve] (finish) {Проверить ограничения, единицы и ответ};
\draw[arr] (start)--(scheme);\draw[arr] (scheme)--(units);\draw[arr] (units)--(table);\draw[arr] (table)--(choice);
\draw[arr] (choice.south west) -| (road.north);
\draw[arr] (choice.south east) -| (times.north);
\draw[arr] (road.south) |- ($(equation.north)+(0,5mm)$) -- (equation.north);
\draw[arr] (times.south) |- ($(equation.north)+(0,5mm)$);
\draw[arr] (equation)--(solve);\draw[arr] (solve)--(finish);
\end{tikzpicture}
\end{center}
\clearpage

\begin{center}{\LARGE\color{darkaccent}\bfseries Тренировочный блок}\par\smallnote{Путь А · 10 заданий · от базовых навыков к задачам уровня ЕГЭ}\end{center}
\vspace{5mm}
\begin{enumerate}
\item \textbf{Перевод единиц.} Переведите: а) $450$ м в километры; б) $18$ мин в часы; в) $0{,}35$ ч в минуты; г) $1$ ч $12$ мин в часы.
\item \textbf{Формула пути.} Автобус двигался со скоростью $72$ км/ч в течение $2$ ч $15$ мин. Какое расстояние он прошёл?
\item \textbf{Догоняющее движение.} Легковой автомобиль выехал вслед за грузовиком, находившимся на $125$ км впереди. Скорость автомобиля на $25$ км/ч больше скорости грузовика. Через сколько часов автомобиль догонит грузовик, если оба продолжают двигаться равномерно?
\item \textbf{Два пешехода.} Из одного места одновременно в одном направлении вышли два пешехода. Скорость первого на $1{,}5$ км/ч больше скорости второго. Через сколько минут расстояние между ними станет равным $450$ м?
\item \textbf{Таблица и уравнение.} Автомобиль выехал из пункта $A$, грузовик --- из пункта $B$, расположенного на том же шоссе на $144$ км впереди пункта $A$. Они выехали одновременно в одном направлении; автомобиль ехал на $18$ км/ч быстрее. Составьте таблицу $v$--$t$--$S$, уравнение и найдите время до встречи.
\item \textbf{Поворот обратно.} Два пешехода одновременно вышли из одного пункта к опушке, находящейся в $8$ км. Скорости --- $7$ и $3$ км/ч. Быстрый дошёл до опушки, повернул назад и встретил медленного. Сколько метров оставалось медленному до опушки?
\item \textbf{Тот же маршрут --- другие числа.} Два туриста вышли одновременно к привалу на расстоянии $5$ км. Скорости --- $6$ и $4$ км/ч. Более быстрый дошёл до привала и пошёл обратно навстречу второму. На каком расстоянии от привала они встретились? Ответ дайте в метрах.
\item \textbf{Разница во времени.} Автобус и грузовик проехали по $24$ км. Грузовик двигался на $5$ км/ч медленнее автобуса, но выехал на $24$ мин раньше, поэтому прибыли они одновременно. Найдите скорость автобуса.
\item \textbf{Квадратное уравнение из задачи.} На участке $36$ км грузовик двигался на $6$ км/ч медленнее автобуса. Он выехал на $30$ мин раньше и прибыл вместе с автобусом. Найдите скорость автобуса. Проверьте ограничения и посторонние корни.
\item \textbf{Поворот и задержка старта.} От старта до опушки $8$ км. Медленный турист начал путь со скоростью $4$ км/ч. Через $15$ мин вслед за ним стартовал быстрый турист со скоростью $8$ км/ч, дошёл до опушки и повернул назад. Они встретились, пока медленный ещё шёл к опушке. На каком расстоянии от опушки произошла встреча? Ответ дайте в километрах.
\end{enumerate}
\clearpage
\begin{center}{\LARGE\color{darkaccent}\bfseries Ответы к тренировочному блоку}\end{center}
\vspace{9mm}
\renewcommand{\arraystretch}{1.65}
\begin{tabularx}{\linewidth}{@{}>{\bfseries}p{.08\linewidth}X@{}}
\toprule
№ & Ответ\\
\midrule
1 & а) $0{,}45$ км; б) $0{,}3$ ч; в) $21$ мин; г) $1{,}2$ ч.\\
2 & $162$ км.\\
3 & $5$ ч.\\
4 & $18$ мин.\\
5 & $8$ ч; возможное уравнение: $x(y+18)=144+xy$.\\
6 & $3200$ м.\\
7 & $1000$ м.\\
8 & $20$ км/ч; допустимо $x>5$.\\
9 & $24$ км/ч; допустимо $x>6$.\\
10 & $2$ км.\\
\bottomrule
\end{tabularx}
\vspace{11mm}
\topic{Краткая памятка}
\begin{tabularx}{\linewidth}{@{}p{.29\linewidth}X@{}}
\toprule
Сигнал в условии & Полезное действие\\\midrule
«догнал», «отставал» & Связать пути через начальную разницу расстояний.\\
«разошлись на 200 м» & Перевести метры в километры и взять разность путей.\\
«дошёл и повернул назад» & Записать $L+x$ для вернувшегося и $L-x$ для встречного.\\
«выехал на 8 минут раньше» & Перевести минуты в часы и составить разность времён.\\
«две дроби в уравнении» & Записать общий знаменатель, проверить его ненулевое значение.\\
«два корня» & Сравнить каждый корень с ограничениями и смыслом задачи.\\
\bottomrule
\end{tabularx}
\vfill
\begin{center}{\footnotesize\color{grayink}Лёвин Артём Александрович эксклюзивно для Анны Банновой}\end{center}
\end{document}
